\chapter{XML Basics}

\section{What is XML?}
XML is the \emph{Extensible Markup Language}. It is designed to
improve the functionality of the Web by providing more flexible and
adaptable information identification.

It is called extensible because it is not a fixed format like HTML (a single,
predefined markup language). Instead, XML is actually a \emph{metalanguage}
-- a language for describing other languages -- which lets you
design your own customized markup languages for limitless different
types of documents. XML can do this because it's written in \emph{SGML}, the
international standard metalanguage for text markup systems (ISO 8879).



\section{XML: Data Elements}
In XML a single data item is enclosed by two tags:
\begin{xmllisting}
  <node_name>ground</node_name>
\end{xmllisting}
an \emph{opening tag} and a \emph{closing tag}. The whole assembly is called an
element.
In XML, an element can have

\begin{center}
  \begin{tabular}{ll}
    \hline\noalign{\smallskip}
    a value    & \xml{<node> ground </node>}\\
    \noalign{\smallskip}\hline\noalign{\smallskip} 
    attributes & \xml{<node name="ground"></node>}\\
    \noalign{\smallskip}\hline\noalign{\smallskip}
    children &  
    \begin{minipage}{10cm}
\begin{xmllisting}
<nodelist>
  <node1>....</node1>
  <node2>....</node2>
</nodelist>
\end{xmllisting}
    \end{minipage}\\
    \noalign{\smallskip}\hline
  \end{tabular}
\end{center}


or a combination of the above.



\section{XML: Empty Elements}
%
%
The XML specification allows for empty elements, i.e.~elements that contain
neither a value nor any childred. Empty, however, does not necessarily mean
that the element does not carry any information, but its name, since an empty
element may still possess attributes. Indeed many of the elements used in the
CFS++ parameter file are empty with information encoded as attribute values.
Note also that, if an element is empty,
\begin{xmllisting}
   <problem name="demo3d"></problem>
\end{xmllisting}
then the opening and the closing tag can, for brevity, be combined into one:
\begin{xmllisting}
   <problem name="demo3d"/>
\end{xmllisting}

